% =====================================================================
\section{Reflection on the use of AI tools}\label{sec:ai}
We used generative AI throughout the project. ChatGPT was used mainly during modelling, for Python code generation, debugging, alternative implementations, review of the geodesic calculations, visualisation design, evaluation and bootstrap code, and interpretation. Claude Code (Anthropic) was used during report preparation, to draft the report from the final notebook, write the figure-regeneration script and cross-check reported numbers against the saved outputs. All modelling decisions and every reported number were checked against the code and its actual outputs. Four concrete examples show where AI helped and where it was wrong.

\textbf{1. Trajectory visualisation (helped after iteration).} The first AI-suggested design overlaid the observed track and all model predictions on one map. With several models the lines were indistinguishable. Over several iterations we moved to one panel per model, fixed geographic extents per drifter so that panels are comparable, solid observed and dashed predicted lines, and in-panel window medians (\cref{fig:traj}). The final design came from judging the plots ourselves, not from the first suggestion.

\textbf{2. Geodesic consistency (AI code was wrong).} AI-assisted code computed observed displacements with a local Cartesian approximation but mapped predictions back with the spherical forward geodesic, mixing two conventions without flagging it. We identified the inconsistency, quantified it (P99 round-trip error \SI{1.54}{km}, maximum \SI{2.98}{km}) and replaced the targets with an inverse-geodesic formulation before refitting all models (\cref{sec:framework}). The lesson is that generated code can be internally plausible yet mathematically inconsistent.

\textbf{3. Methodological wording (needed correction).} AI-generated descriptions did not reliably distinguish our single drifter-disjoint hold-out from $K$-fold cross-validation, or rolling 24-hour forecasts re-initialised from observations from recursive trajectory forecasting. We checked the wording against the code, which uses \texttt{GroupShuffleSplit} and no $K$-fold routine, and corrected it, because these distinctions determine what the results can claim.

\textbf{4. Report preparation (helped, with one error caught).} When drafting the bibliography, the assistant proposed a DOI from memory for a Gulf Stream reference. It resolved to an unrelated article and was discarded once every DOI was checked against Crossref. On the positive side, systematic cross-checking surfaced issues we had not documented: the undrogued share of the cohort, a shared legend in an existing figure whose error values apply to only one panel, and the default early-stopping behaviour of the boosting estimator. We verified each against the data or library documentation before including it.

Overall, AI accelerated coding, plotting and drafting. It was least reliable where correctness depended on domain conventions or precise experimental definitions, and human verification against outputs remained essential.
